\documentclass[11pt,a4paper]{article}
\usepackage[T1]{fontenc}
\usepackage[utf8]{inputenc}
\usepackage{lmodern}
\usepackage{amsmath,amssymb,amsthm,mathtools}
\usepackage{geometry}
\usepackage{hyperref}
\geometry{margin=2.5cm}

\title{TGP: Algebraiczne wyprowadzenia dla szczególnych konfiguracji \texorpdfstring{$N$}{N}-ciał w sektorze pairwise}
\author{}
\date{}

\theoremstyle{plain}
\newtheorem{lemma}{Lemma}
\newtheorem{proposition}{Propozycja}
\theoremstyle{remark}
\newtheorem*{remark}{Uwaga}

\begin{document}
\maketitle

\section{Cel notatki}

Celem tej notatki jest zapisanie w zwartej formie wyprowadzeń dla klas szczególnych konfiguracji w sektorze \emph{pairwise} TGP, dla których warunek równowagi redukuje się do pojedynczego równania algebraicznego drugiego stopnia dla parametru skali.

Punktem wyjścia jest potencjał pairwise dla równych źródeł $C$:
\begin{equation}
V_2(d) = -\frac{4\pi C^2}{d} + \frac{8\pi \beta C^2}{d^2} - \frac{24\pi \gamma C^3}{d^3}.
\end{equation}

W wielu wcześniejszych rozważaniach szczególny przypadek odpowiadał $\gamma = \beta$, ale poniżej pozostawiamy oba parametry jawnie, ponieważ algebraiczne wyprowadzenia nie wymagają ich utożsamienia.

\section{Konfiguracje jednoskalowe}

Rozważmy układ $N$ równych źródeł, którego geometria ma stały kształt i zależy wyłącznie od jednej skali $\lambda$. Oznacza to, że wszystkie odległości między ciałami mają postać
\begin{equation}
d_{ij} = a_{ij} \lambda,
\end{equation}
 gdzie $a_{ij} > 0$ są bezwymiarowymi stałymi geometrycznymi określającymi kształt konfiguracji.

Całkowita energia pairwise wynosi wtedy
\begin{equation}
V(\lambda) = \sum_{i<j} V_2(d_{ij})
= \sum_{i<j} \left(
-\frac{4\pi C^2}{a_{ij}\lambda}
+ \frac{8\pi\beta C^2}{a_{ij}^2\lambda^2}
- \frac{24\pi\gamma C^3}{a_{ij}^3\lambda^3}
\right).
\end{equation}

Wprowadzamy sumy geometryczne:
\begin{equation}
S_1 = \sum_{i<j}\frac{1}{a_{ij}},
\qquad
S_2 = \sum_{i<j}\frac{1}{a_{ij}^2},
\qquad
S_3 = \sum_{i<j}\frac{1}{a_{ij}^3}.
\end{equation}

Wtedy energia przyjmuje postać
\begin{equation}
V(\lambda) = -\frac{4\pi C^2}{\lambda}S_1
+ \frac{8\pi\beta C^2}{\lambda^2}S_2
- \frac{24\pi\gamma C^3}{\lambda^3}S_3.
\end{equation}

\begin{lemma}[Ogólna redukcja algebraiczna dla konfiguracji jednoskalowych]
Dla dowolnej konfiguracji o stałym kształcie i jednej skali $\lambda$, warunek równowagi w sektorze pairwise,
\begin{equation}
\frac{dV}{d\lambda}=0,
\end{equation}
redukuje się do równania kwadratowego
\begin{equation}
\boxed{S_1\lambda^2 - 4\beta S_2\lambda + 18\gamma C S_3 = 0.}
\end{equation}
\end{lemma}

\begin{proof}
Różniczkujemy energię po $\lambda$:
\begin{equation}
\frac{dV}{d\lambda}
= \frac{4\pi C^2}{\lambda^2}S_1
- \frac{16\pi\beta C^2}{\lambda^3}S_2
+ \frac{72\pi\gamma C^3}{\lambda^4}S_3.
\end{equation}

Warunek stacjonarności ma więc postać
\begin{equation}
\frac{4\pi C^2}{\lambda^2}S_1
- \frac{16\pi\beta C^2}{\lambda^3}S_2
+ \frac{72\pi\gamma C^3}{\lambda^4}S_3 = 0.
\end{equation}

Po pomnożeniu przez $\lambda^4/(4\pi C^2)$ dostajemy
\begin{equation}
S_1\lambda^2 - 4\beta S_2\lambda + 18\gamma C S_3 = 0.
\end{equation}
To kończy dowód.
\end{proof}

\begin{remark}
Powyższy wynik nie rozwiązuje ogólnego problemu $N$-ciał. Pokazuje natomiast, że szeroka klasa konfiguracji samopodobnych w sektorze pairwise redukuje się do pojedynczego równania algebraicznego dla parametru skali.
\end{remark}

\section{Przykłady}

\subsection{Dwa ciała}

Dla dwóch ciał istnieje tylko jedna odległość, więc można przyjąć
\begin{equation}
d = \lambda, \qquad a_{12}=1.
\end{equation}
Stąd
\begin{equation}
S_1=S_2=S_3=1.
\end{equation}
Warunek równowagi daje
\begin{equation}
\boxed{\lambda^2 - 4\beta \lambda + 18\gamma C = 0.}
\end{equation}

W przypadku $\gamma = \beta$ otrzymujemy
\begin{equation}
\lambda^2 - 4\beta \lambda + 18\beta C = 0.
\end{equation}

\subsection{Trójkąt równoboczny}

Niech $d$ oznacza długość boku trójkąta równobocznego. Wszystkie trzy odległości są równe $d$, zatem
\begin{equation}
S_1=S_2=S_3=3.
\end{equation}
Warunek równowagi ma postać
\begin{equation}
3d^2 - 12\beta d + 54\gamma C = 0,
\end{equation}
co po skróceniu przez $3$ daje
\begin{equation}
\boxed{d^2 - 4\beta d + 18\gamma C = 0.}
\end{equation}

W szczególności dla $\gamma = \beta$:
\begin{equation}
d^2 - 4\beta d + 18\beta C = 0.
\end{equation}

\subsection{Tetraedr foremny}

Niech $a$ oznacza długość krawędzi tetraedru foremnego. Liczba par wynosi $6$, a wszystkie odległości są równe $a$. Zatem
\begin{equation}
S_1=S_2=S_3=6.
\end{equation}
Stąd
\begin{equation}
6a^2 - 24\beta a + 108\gamma C = 0,
\end{equation}
więc po skróceniu:
\begin{equation}
\boxed{a^2 - 4\beta a + 18\gamma C = 0.}
\end{equation}

Widać więc, że dwa ciała, trójkąt równoboczny i tetraedr foremny dają to samo równanie dla wspólnej odległości, ponieważ wszystkie pary są równoodległe.

\subsection{Kwadrat}

Niech $a$ oznacza długość boku kwadratu. Występują następujące odległości:
\begin{itemize}
    \item $4$ boki długości $a$,
    \item $2$ przekątne długości $\sqrt{2}\,a$.
\end{itemize}

Wobec tego
\begin{equation}
S_1 = 4 + \frac{2}{\sqrt{2}} = 4 + \sqrt{2},
\end{equation}
\begin{equation}
S_2 = 4 + \frac{2}{2} = 5,
\end{equation}
\begin{equation}
S_3 = 4 + \frac{2}{2\sqrt{2}} = 4 + \frac{1}{\sqrt{2}}.
\end{equation}

Warunek równowagi przyjmuje postać
\begin{equation}
\boxed{(4+\sqrt{2})a^2 - 20\beta a + 18\gamma C\left(4+\frac{1}{\sqrt{2}}\right)=0.}
\end{equation}

Jest to pierwszy nietrywialny przykład, w którym nie wszystkie odległości są jednakowe, a mimo to całość redukuje się do pojedynczego równania kwadratowego.

\subsection{Symetryczny układ współliniowy trzech ciał}

Rozważmy pozycje
\begin{equation}
(-x,0,+x).
\end{equation}
Odległości między ciałami są równe
\begin{equation}
x,\quad x,\quad 2x.
\end{equation}

Stąd
\begin{equation}
S_1 = 2 + \frac{1}{2} = \frac{5}{2},
\end{equation}
\begin{equation}
S_2 = 2 + \frac{1}{4} = \frac{9}{4},
\end{equation}
\begin{equation}
S_3 = 2 + \frac{1}{8} = \frac{17}{8}.
\end{equation}

Warunek równowagi daje więc
\begin{equation}
\frac{5}{2}x^2 - 4\beta\frac{9}{4}x + 18\gamma C\frac{17}{8}=0,
\end{equation}
czyli po uproszczeniu
\begin{equation}
\boxed{10x^2 - 36\beta x + 153\gamma C = 0.}
\end{equation}

\subsection{Regularny \texorpdfstring{$n$}{n}-kąt}

Niech $R$ oznacza promień okręgu opisanego na regularnym $n$-kącie. Odległości od ustalonego wierzchołka do pozostałych mają postać
\begin{equation}
d_k = 2R\sin\frac{\pi k}{n},
\qquad k=1,\dots,n-1.
\end{equation}

Wygodnie wprowadzić sumy geometryczne zależne od $n$:
\begin{equation}
\Sigma_p(n) = \sum_{i<j}\frac{1}{a_{ij}^p},
\end{equation}
 gdzie $a_{ij}$ są bezwymiarowymi odległościami w regularnym $n$-kącie po wyłączeniu skali $R$.

Równoważnie można je zapisać przez klasy odległości cięciwowych. Wtedy warunek równowagi ma postać
\begin{equation}
\boxed{\Sigma_1(n)R^2 - 4\beta\Sigma_2(n)R + 18\gamma C\Sigma_3(n)=0.}
\end{equation}

Dla $n=3$ otrzymujemy trójkąt równoboczny jako przypadek szczególny.

\begin{remark}
Dokładna postać sum $\Sigma_p(n)$ zależy od przyjętej konwencji liczenia par i może być zapisana albo jako suma po wszystkich parach $(i,j)$, albo jako suma po klasach odległości w regularnym $n$-kącie z odpowiednimi krotnościami.
\end{remark}

\subsection{Równoodległy układ współliniowy \texorpdfstring{$N$}{N}-ciał}

Rozważmy $N$ równych ciał na prostej, rozmieszczonych co odległość $x$:
\begin{equation}
0,\ x,\ 2x,\ \dots,\ (N-1)x.
\end{equation}

Odległości między parami mają postać
\begin{equation}
mx,\qquad m=1,\dots,N-1,
\end{equation}
przy czym każda odległość $mx$ występuje z krotnością $(N-m)$.

Dlatego sumy geometryczne wynoszą
\begin{equation}
S_p(N)=\sum_{m=1}^{N-1}\frac{N-m}{m^p},
\qquad p=1,2,3.
\end{equation}

Warunek równowagi redukuje się do
\begin{equation}
\boxed{S_1(N)x^2 - 4\beta S_2(N)x + 18\gamma C S_3(N)=0.}
\end{equation}

Dla $N=3$ otrzymujemy ponownie wynik dla układu $(-x,0,+x)$.

\section{Interpretacja wyniku}

Najmocniejszy wniosek z powyższych obliczeń nie brzmi: ``problem $N$-ciał stał się algebraicznie rozwiązywalny''.

Wniosek jest słabszy, ale nadal istotny:
\begin{quote}
W sektorze pairwise TGP szeroka klasa konfiguracji o stałym kształcie i jednej skali redukuje warunek równowagi do pojedynczego równania kwadratowego dla parametru skali.
\end{quote}

To oznacza, że pairwise TGP nie daje pełnego rozwiązania ogólnego problemu $N$-ciał, ale selekcjonuje klasy figur, dla których struktura energii jest wyjątkowo prosta.

\section{Granice zastosowania}

Wszystkie powyższe wyprowadzenia dotyczą wyłącznie sektora \emph{pairwise}. W pełnej teorii energia ma schematycznie postać
\begin{equation}
V_{\mathrm{total}} = \sum_{i<j}V_2(d_{ij}) + \sum_{i<j<k}V_3(d_{ij},d_{ik},d_{jk}),
\end{equation}
co oznacza, że po uwzględnieniu rzeczywistego członu trójciałowego prostota algebraiczna może zostać utracona.

W szczególności, jeśli człon $V_3$ zawiera nietrywialne całki overlapu albo screening typu Yukawy, to pełna energia może przestać być prostą kombinacją składników $1/\lambda$, $1/\lambda^2$, $1/\lambda^3$, a wtedy redukcja do równania kwadratowego nie musi już zachodzić.

\section{Dalsze kierunki}

Naturalne dalsze kroki:
\begin{enumerate}
    \item sprawdzić, które z powyższych rodzin przeżywają po dołączeniu członu 3-body jako perturbacji,
    \item zbadać stabilność małych zaburzeń wokół uzyskanych konfiguracji,
    \item sklasyfikować wszystkie figury jednoskalowe, dla których sumy geometryczne $S_1,S_2,S_3$ przyjmują szczególnie proste postaci,
    \item sprawdzić, czy istnieją rodziny samopodobnych konfiguracji w 3D (np. inne wielościany foremne lub półforemne), które również poddają się tej redukcji.
\end{enumerate}

\end{document}
